\documentclass[12pt,a4paper]{report}
\usepackage[utf8]{inputenc}
\usepackage{amsmath}
\usepackage{amsfonts}
\usepackage{amssymb}
\usepackage{amsthm}
\usepackage{hyperref}

\usepackage{multicol}
\usepackage{fancyhdr}
\usepackage[inline]{enumitem}
\usepackage{tikz}
\usepackage{tikz-cd}
\usetikzlibrary{calc}
\usetikzlibrary{shapes.geometric}
\usetikzlibrary{positioning}
\usepackage[margin=0.5in]{geometry}
\usepackage{xcolor}

\hypersetup{
    colorlinks=true,
    linkcolor=blue,
    filecolor=magenta,      
    urlcolor=cyan,
    pdftitle={Tensors},
    pdfpagemode=FullScreen,
    }

%\urlstyle{same}

\newcommand{\CLASSNAME}{Math 5050 -- Special Topics: Differential Geometry}
\newcommand{\STUDENTNAME}{Paul Carmody}
\newcommand{\ASSIGNMENT}{Section 24: de Rhan Cohomology }
\newcommand{\DUEDATE}{November 26, 2025}
\newcommand{\PROFESSOR}{Professor Berchenko-Kogan}
\newcommand{\SEMESTER}{Fall 2025}
\newcommand{\SCHEDULE}{TBD}
\newcommand{\ROOM}{Remote}

\newcommand{\MMN}{M_{m\times n}}
\newcommand{\FF}{\mathcal{F}}

\pagestyle{fancy}
\fancyhf{}
\chead{ \fancyplain{}{\CLASSNAME} }
%\chead{ \fancyplain{}{\STUDENTNAME} }
\rhead{\thepage}
\input{../MyMacros}

\newcommand{\RED}[1]{\textcolor{red}{#1}}
\newcommand{\BLUE}[1]{\textcolor{blue}{#1}}

\newcommand{\SUPP}{\operatorname{supp}}
\newcommand{\CLOSURE}{\operatorname{closure of}}
\newcommand{\BD}{\operatorname{bd}}
\newcommand{\HHN}{\mathcal{H}^n}

\begin{document}

\begin{center}
	\Large{\CLASSNAME -- \SEMESTER} \\
	\large{ w/\PROFESSOR}
\end{center}
\begin{center}
	\STUDENTNAME \\
	\ASSIGNMENT -- \DUEDATE\\
\end{center} 

\textbf{Key Points}

\HLINE
\begin{enumerate}[label=24.\arabic*]

\item \textbf{Nowhere vanishing $1$-forms}

Prove that a nowhere vanishing $1$-form on a compact manifold cannot exist.

\BLUE{Not True: On the circle $S^1$ the form $d\theta$ (in the angular coordinate $\theta$) is smooth 1-form that never vanishes.  $S^1$ is compact so compactness alone does not forbid nowhere-vanishing 1-forms.
}

\item \textbf{Cohomology in degree zero}

Suppose a manifold $M$ has infinitely many connected components.  Compute its de Rham Cohomology vector space $H^0(M)$ in degree 0. (\textit{Hint:}  by second countability,  the number of connect components of a manifold is countable.)

\BLUE{Let $\dim M = n$ with $n$ connected components.  If $\omega = (\omega_1,\dots,\omega_n)$ is closed but not exact, then, $d\omega = 0$.  Since, 
\begin{align*}
	\omega &= \sum_{i=1}^n \omega_i x^i\\
	d\omega &= \sum_{i=1}^n d\omega_i \,dx^i \\
	\therefore d\omega_i &= 0, \, \forall i=1,\dots,n
\end{align*}and $\omega_i$ are constants for all $i=1,\dots,n$.  Now if $\dim M = \infty$, we recall that $M$ must be second countable.  Therefore, we can have $\omega = (\omega_1, \dots, )$ which allows a sum to infinity.
}

\end{enumerate}
\end{document}
